\documentclass[11pt,a4paper]{article}

\usepackage[T1]{fontenc}
\usepackage[utf8]{inputenc}
\usepackage[ngerman]{babel}
\usepackage{amsmath,amssymb,amsfonts,amsthm}
\usepackage{geometry}
\usepackage{hyperref}
\usepackage{graphicx}
\usepackage{enumitem}

\geometry{margin=2.5cm}

\title{
	Rayleighs Fluch als Projektionsphänomen:\\
	Eine EABC-Interpretation moderner Quantenmetrologie
}

\author{
	Thomas Hoffbauer
}

\date{\today}

\newtheorem{definition}{Definition}
\newtheorem{satz}{Satz}
\newtheorem{lemma}{Lemma}
\newtheorem{korollar}{Korollar}
\theoremstyle{remark}
\newtheorem{bemerkung}{Bemerkung}

\begin{document}
	
	\maketitle
	
	\begin{abstract}
		
		Jüngste Experimente aus der Quantenoptik zeigen, dass die klassische
		Rayleigh-Grenze nicht durch einen fundamentalen Informationsverlust
		verursacht wird, sondern durch die Wahl einer ungeeigneten Messbasis.
		Durch die Projektion optischer Zustände auf geeignete Moden konnte die
		Messpräzision um mehr als eine Größenordnung gesteigert werden.
		In der vorliegenden Arbeit wird vorgeschlagen, diese Beobachtung als
		Manifestation eines allgemeineren informationsgeometrischen Prinzips zu
		interpretieren.
		
		Innerhalb des EABC-Modells wird gezeigt, dass scheinbare
		Informationsgrenzen als Projektionsartefakte verstanden werden können.
		Information geht nicht verloren, sondern wird lediglich in einer
		ungeeigneten Basis dargestellt.
		Die Analyse legt nahe, dass Quantenmetrologie, Modenfilterung und die
		EABC-Struktur verschiedene Erscheinungsformen eines gemeinsamen
		strukturellen Prinzips sind.
		
	\end{abstract}
	
	\section{Einleitung}
	
	Seit der Arbeit von Lord Rayleigh gilt die Auflösungsgrenze optischer
	Systeme als fundamentale Einschränkung der Messgenauigkeit.
	
	Werden zwei Lichtquellen oder Frequenzen ausreichend nahe
	zusammengebracht, verschmelzen ihre Signaturen scheinbar zu einem
	einzigen Signal.
	
	Die klassische Interpretation lautet:
	
	\begin{equation}
		\text{kleiner Abstand}
		\quad\Longrightarrow\quad
		\text{Informationsverlust}.
	\end{equation}
	
	Moderne Arbeiten aus der Quantenmetrologie zeigen jedoch ein anderes
	Bild.
	
	Die Information bleibt vollständig im Zustand erhalten.
	Verloren geht lediglich die Fähigkeit einer gegebenen Messbasis,
	diese Information sichtbar zu machen.
	
	Diese Arbeit untersucht die Konsequenzen dieser Sichtweise im Rahmen des
	EABC-Modells.
	
	\section{Rayleighs Fluch}
	
	Ein optischer Zustand kann allgemein als Modenüberlagerung geschrieben
	werden:
	
	\begin{equation}
		|\psi\rangle
		=
		\sum_i c_i |u_i\rangle.
	\end{equation}
	
	Hierbei bilden
	
	\[
	|u_i\rangle
	\]
	
	eine vollständige Modenbasis.
	
	Konventionelle Messungen erfassen typischerweise nur die
	Intensitätsverteilung
	
	\begin{equation}
		I(x)
		=
		|\psi(x)|^2.
	\end{equation}
	
	Dadurch werden zahlreiche Freiheitsgrade des Zustands ignoriert.
	
	Die moderne Quantenmetrologie verwendet stattdessen gezielte
	Projektionen
	
	\begin{equation}
		M_i
		=
		\langle u_i|\psi\rangle,
	\end{equation}
	
	wodurch verborgene Informationsanteile zugänglich werden.
	
	\section{EABC-Zustandsräume}
	
	Das EABC-Modell beschreibt physikalische Zustände nicht als isolierte
	Objekte, sondern als Elemente eines strukturierten Zustandsraums.
	
	Ein Zustand besitzt die Darstellung
	
	\begin{equation}
		|\Psi\rangle
		=
		\sum_{\alpha}
		c_{\alpha}
		|E_{\alpha}\rangle.
	\end{equation}
	
	Die Basisvektoren
	
	\[
	|E_{\alpha}\rangle
	\]
	
	repräsentieren EABC-Moden.
	
	Diese können beispielsweise sein:
	
	\begin{itemize}
		\item Schalenzustände,
		\item Nachbarschaftsklassen,
		\item Holonomieklassen,
		\item diskrete Weyl-Moden,
		\item EABC-Symmetrieklassen.
	\end{itemize}
	
	Die Beobachtung erfolgt über Projektionsoperatoren
	
	\begin{equation}
		P_{\alpha}
		=
		|E_{\alpha}\rangle
		\langle E_{\alpha}|.
	\end{equation}
	
	Die beobachtbare Information ergibt sich dann zu
	
	\begin{equation}
		I_{\alpha}
		=
		\langle \Psi|
		P_{\alpha}
		|\Psi\rangle.
	\end{equation}
	
	\section{Informationsverlust als Projektionsfehler}
	
	Die zentrale These lautet:
	
	\begin{satz}[Projektionsprinzip]
		Sei
		
		\[
		|\Psi\rangle
		\in \mathcal H
		\]
		
		ein Zustand eines strukturierten Zustandsraums.
		
		Eine beobachtete Informationsgrenze kann entstehen, obwohl die
		Information vollständig im Zustand enthalten ist.
		
		In diesem Fall beruht die Grenze auf einer ungeeigneten Wahl der
		Projektionsbasis.
	\end{satz}
	
	\begin{proof}
		
		Sei
		
		\[
		\mathcal B
		=
		\{|b_i\rangle\}
		\]
		
		eine Beobachtungsbasis.
		
		Die beobachtbare Information besitzt die Form
		
		\[
		I_i
		=
		|\langle b_i|\Psi\rangle|^2.
		\]
		
		Ist die Basis nicht an die tatsächliche Struktur des Zustands angepasst,
		so können mehrere unterschiedliche Zustände identische
		Projektionsbilder erzeugen.
		
		Wird dagegen eine angepasste Basis
		
		\[
		\mathcal E
		=
		\{|E_\alpha\rangle\}
		\]
		
		gewählt, können dieselben Zustände unterscheidbar werden.
		
		Die Informationsgrenze verschwindet somit ohne Änderung des
		physikalischen Zustands.
		
		Folglich liegt kein Informationsverlust vor, sondern ein
		Projektionsartefakt.
	\end{proof}
	
	\section{EABC-Interpretation der Oxford-Experimente}
	
	Die jüngsten Experimente mit Caesium-Quantenspeichern zeigen, dass eine
	geeignete Modenprojektion die Messgenauigkeit erheblich steigern kann.
	
	Innerhalb der EABC-Sichtweise bedeutet dies:
	
	\begin{equation}
		\text{Messinformation}
		\neq
		\text{Signalstärke},
	\end{equation}
	
	sondern
	
	\begin{equation}
		\text{Messinformation}
		=
		\text{Strukturinformation}.
	\end{equation}
	
	Die Aufgabe des Quantenspeichers besteht darin, eine bessere
	Darstellung des Zustandsraums bereitzustellen.
	
	Damit wird die Information nicht erzeugt, sondern sichtbar gemacht.
	
	\section{Informationsgeometrische Formulierung}
	
	Die beobachtbare Realität hängt sowohl vom Zustand als auch von der
	gewählten Projektion ab.
	
	Formal schreiben wir
	
	\begin{equation}
		R
		=
		\mathcal P(\Psi),
	\end{equation}
	
	wobei
	
	\[
	\mathcal P
	\]
	
	den Beobachtungsoperator bezeichnet.
	
	Zwei Beobachter können daher dieselbe physikalische Struktur
	unterschiedlich wahrnehmen:
	
	\begin{equation}
		\mathcal P_1(\Psi)
		\neq
		\mathcal P_2(\Psi).
	\end{equation}
	
	Die Differenz liegt nicht im Zustand, sondern in der
	Informationsgeometrie der Beobachtung.
	
	\section{Konsequenzen}
	
	Die hier vorgestellte Interpretation legt eine allgemeine Sichtweise
	nahe:
	
	\begin{enumerate}[label=\arabic*.]
		
		\item Rayleigh-Grenzen sind nicht notwendigerweise fundamental.
		
		\item Informationsverlust ist häufig ein Projektionsproblem.
		
		\item EABC-Strukturen können als natürliche Modenräume interpretiert
		werden.
		
		\item Quantenmetrologie und EABC besitzen eine gemeinsame
		informationsgeometrische Grundlage.
		
		\item Die Wahl der Beobachtungsbasis wird zu einem physikalischen
		Freiheitsgrad.
		
	\end{enumerate}
	
	\section{Ausblick}
	
	Die vorliegende Arbeit versteht sich als konzeptioneller Beitrag zur
	Verbindung von Quantenmetrologie und EABC-Theorie.
	
	Zukünftige Untersuchungen könnten konkrete EABC-Basen konstruieren und
	deren Informationsgehalt mit Hilfe von Fisher-Information,
	Quanten-Fisher-Information und Holonomieoperatoren analysieren.
	
	Von besonderem Interesse ist die Frage, ob bekannte
	Auflösungsgrenzen anderer physikalischer Systeme ebenfalls als
	Projektionsartefakte interpretiert werden können.
	
	\section*{Danksagung}
	
	Der Autor dankt den Forschern der modernen Quantenmetrologie für die
	experimentellen Ergebnisse, die den Anlass zu den hier entwickelten
	Überlegungen gaben.
	
	\begin{thebibliography}{9}
		
		\bibitem{rayleigh}
		J. W. Strutt (Lord Rayleigh),
		\emph{Investigations in Optics},
		Philosophical Magazine,
		1879.
		
		\bibitem{tsang}
		M. Tsang, R. Nair, X.-M. Lu,
		Quantum Theory of Superresolution,
		Physical Review X,
		2016.
		
		\bibitem{oxford}
		Oxford University,
		Quantum Memory Sharpens Optical Measurements Beyond Rayleigh's Curse,
		Nature Sensors,
		2026.
		
		\bibitem{helstrom}
		C. Helstrom,
		\emph{Quantum Detection and Estimation Theory},
		Academic Press,
		1976.
		
	\end{thebibliography}
	
\end{document}